\section{II. QUANTUM CIRCUIT}

\subsection{Paragraf 1: Implementasi Sirkuit}
\begin{frame}{Paragraf 1 (1-8)}
    \begin{itemize}
        \item \textbf{1. model consideration:} Syarat matriks non-negatif $\sigma_N$ dengan \textit{trace} satu.
        \item \textbf{2. asumsi unitary:} Efisiensi pembangkitan uniter $e^{it\sigma_N}$ sebagai prasyarat.
        \item \textbf{4. kasus paper:} Akses ke banyak salinan \textit{state} memungkinkan uniter efisien meski tidak \textit{sparse}.
        \item \textbf{5. best way:} Memilih metode Lloyd et al. (Ref [14]) untuk akurasi terkendali.
        \item \textbf{7. tujuan algoritma:} Penentuan $r$ \textit{eigenvalue} terbesar dan \textit{eigenvector} terkait.
        \item \textbf{8. ilustrasi:} Referensi Gambar 1 untuk dekomposisi gerbang $n + \log N$ qubit.
    \end{itemize}
\end{frame}

\begin{frame}{Paragraf 1 (9-15)}
    \begin{itemize}
        \item \textbf{9. a priori eigenvectors:} Tidak mengetahui \textit{eigenvector} menghambat QPE langsung.
        \item \textbf{11. random state:} Sistem diinisialisasi pada $\ket{b}$ sebagai konsekuensi stokastik.
        \item \textbf{13. QFT and application:} Peran QFT menciptakan \textit{entanglement} antara estimasi \textit{eigenvalue} dan \textit{eigenvector}.
        \item \textbf{14. $\Lambda_j^{(n)}$:} Notasi representasi biner $n$-bit untuk presisi.
        \item \textbf{15. aproksimasi bagus:} Komponen \textit{eigenvalue} tertinggi mendominasi hasil proyeksi.
    \end{itemize}
\end{frame}

\begin{frame}{Paragraf 1 (16-21)}
    \begin{itemize}
        \item \textbf{16. $\ket{y^{(n)}}$:} Definisi vektor proyeksi untuk mengekstraksi \textit{eigenvector}.
        \item \textbf{18. kondisi NISQ:} Presisi bit rendah menyebabkan ambiguitas identifikasi antar \textit{eigenvector}.
        \item \textbf{19. verifikasi:} Menggunakan berbagai \textit{random state} $\ket{c}$ untuk konsistensi hasil.
        \item \textbf{21. heuristik:} Solusi peningkatan $n$-bit jika identifikasi unik gagal.
    \end{itemize}
\end{frame}

\begin{frame}{Pertanyaan yang dijawab pada Paragraf 1}
    \centering
    \small
    \begin{tabular}{lp{4.5cm}p{6cm}}
        \toprule
        Kalimat ke & Pertanyaan & Jawaban \\
        \midrule
        1 & Batasan teoretis matriks? & Non-negatif $\sigma_N$ dengan $tr[\sigma_N]=1$. \\
        9 & Mengapa QPE tidak langsung? & \textit{Eigenvector} tidak diketahui (\textit{a priori}). \\
        13 & Peran QFT? & Menciptakan \textit{entanglement} estimasi eigen. \\
        18 & Tantangan NISQ? & Presisi rendah memicu ambiguitas eigen. \\
        \bottomrule
    \end{tabular}
\end{frame}

\subsection{Paragraf 2: Perbaikan Sekuensial}
\begin{frame}{Paragraf 2}
    \begin{itemize}
        \item \textbf{1. asumsi n-bit:} Anggapan presisi cukup untuk memisahkan \textit{eigenvector} unik.
        \item \textbf{2. constraint:} Limitasi qubit, derau chip, dan kesalahan operasional.
        \item \textbf{3. mulai protokol:} Inisialisasi $\ket{b_0}$ dan proyeksi ke ruang $\ket{y^{(n)}}$.
        \item \textbf{5. error cancellation:} Feedback hasil proyeksi untuk mekanisme pembatalan kesalahan koheren.
        \item \textbf{6. limitasi:} Titik dekoherensi dan galat statistik pengukuran sebagai batas fidelitas.
        \item \textbf{7. mitigasi:} Pengukuran pada berbagai basis dan rata-rata untuk reduksi galat sistematis.
    \end{itemize}
\end{frame}

\begin{frame}{Pertanyaan yang dijawab pada Paragraf 2}
    \centering
    \small
    \begin{tabular}{lp{4.5cm}p{6cm}}
        \toprule
        Kalimat ke & Pertanyaan & Jawaban \\
        \midrule
        1 & Kapan proses perbaikan mulai? & Saat presisi $n$-bit cukup pisahkan \textit{eigenvector}. \\
        2 & Faktor pembatas perbaikan? & Limitasi qubit, derau chip, galat operasional. \\
        5 & Mengapa reuse hasil lama? & Tingkatkan fidelitas via pembatalan galat koheren. \\
        7 & Reduksi galat sistematis? & Pengukuran berbagai basis dan ambil rata-rata. \\
        \bottomrule
    \end{tabular}
\end{frame}
